\section{PyTorch 编译优化（torch.compile）}

\subsection{自动融合的魔力}

\texttt{torch.compile} 是 PyTorch 2.0 引入的 JIT 编译器，可以\textbf{自动}将朴素 PyTorch 代码优化为融合的 Triton 内核：

\begin{lstlisting}[caption={torch.compile 使用示例}]
# 原始代码（未融合，约 8.1ms）
def manual_gelu(x):
    return 0.5 * x * (1 + torch.tanh(
        0.79788456 * (x + 0.044715 * x * x * x)
    ))

# 一行代码即可优化（约 1.47ms）
compiled_gelu = torch.compile(manual_gelu)
\end{lstlisting}

\texttt{torch.compile} 在底层做了什么？
\begin{enumerate}
    \item 分析 Python 代码的计算图
    \item 识别可融合的操作
    \item 自动生成优化的 Triton 内核代码
    \item 编译并缓存生成的代码
\end{enumerate}

Profiling 证实：编译后的 GeLU 在底层调用了一个名为 \texttt{fused\_add\_multiply\_tanh} 的 Triton 内核——与我们手写的 Triton 代码做了类似的融合，但自动生成的代码略有优化。

\subsection{torch.compile 的适用场景}

\begin{importantbox}{何时用 torch.compile，何时手写 Triton？}
\begin{itemize}
    \item \textbf{简单的算子融合}（如 GeLU、pointwise 操作链）：\texttt{torch.compile} 完全够用，不需要手写内核
    \item \textbf{矩阵乘法调度优化}：\texttt{torch.compile} 可以自动选择最优的内核实现，带来约 10\% 免费加速
    \item \textbf{复杂的算法优化}（如 FlashAttention）：仍需手写 Triton/CUDA——编译器无法自动发现 online softmax 这样的算法创新
    \item \textbf{硬件特定优化}（如 FlashAttention-3 利用 H100 特性）：需要手写底层代码
\end{itemize}
\end{importantbox}

\begin{warningbox}{不要盲目手写 CUDA 内核}
``不要回家后就给语言模型的每个组件都写 CUDA 内核——那可能不是你时间的最佳用途。'' 但如果你在开发新架构，某个复杂组件的 GPU 利用率很低，而你认为可以改善——那就是拿出 Triton 的时候。
\end{warningbox}

\subsection{本章小结}

\texttt{torch.compile} 是``低垂的果实''——一行代码就能获得接近手写内核的性能。它在底层生成 Triton 代码来实现算子融合。对于标准操作，优先使用 \texttt{torch.compile}，只在需要非平凡优化时才手写 Triton 内核。

%% ========== Section 9: Softmax ==========

